%% toc: P15 166 166.1
% Shannon's entropy functional: domain, units, characterization, and what it does not measure.
\chapter{Information and Description-Based Inference}
\label{sec:p15-w02b-p15-040-information}

A probability distribution over a set of possible outcomes can be
summarized by a single number that measures how much choice, or
uncertainty, it leaves about which outcome occurs. Shannon called this
number the entropy of the distribution. This chapter states what that
number is, on exactly which distributions it is defined, in which units
it is measured, and why the logarithmic formula is forced by three
requirements. Later sections of the chapter treat joint and conditional
entropy and the entropy of densities
(\cref{sec:p15-w02b-p15-041}), and maximum-entropy inference under
declared constraints (\cref{sec:p15-w02b-p15-044-maxent-existence}).

\emph{Inputs.} A probability vector on a finite or countably infinite
set of outcomes, and a choice of logarithmic unit. \emph{Conclusions.}
The value of the entropy, possibly \(+\infty\) in the countable case;
its conversion between units; its extreme values; and the theorem that
a measure of uncertainty that is continuous, increases with the number
of equally likely outcomes, and decomposes over successive choices must
be a positive multiple of \(-\sum p_i\log p_i\). The presentation
follows Section~6 and Appendix~2 of Shannon's corrected reprint of the
1948 paper \cite[pp.~10--12, 28--29]{Shannon1948MathematicalTheory},
whose page numbers are those of the reprint.

\section{Entropy of discrete distributions and its characterization}
\label{sec:p15-w02b-p15-040-discrete}

\begin{hbdefinition}{def:p15-w02b-p15-040-entropy}{The discrete entropy functional and its units}
Let \(I\) be a nonempty finite or countably infinite index set. A
\emph{probability vector} on \(I\) is a function
\(p:I\to[0,1]\), \(i\mapsto p_i\), with \(\sum_{i\in I}p_i=1\); write
\(\Delta(I)\) for the set of them and
\(\Delta_n\defeq\Delta(\{1,\ldots,n\})\) for \(n\geq1\). Logarithms
written \(\log\) are natural. Fix a constant \(K>0\) and define
\[
 H(p)\defeq-K\sum_{i\in I}p_i\log p_i
      \defeq K\sum_{i\in I:\,p_i>0}p_i\log\frac1{p_i}
      \in[0,+\infty].
\]
Equivalently, the full sum is used with the convention \(0\log0=0\).
For \(p\in\Delta_n\) write \(H(p_1,\ldots,p_n)\) for \(H(p)\).

For a base \(b>1\), the entropy \emph{in base \(b\)} is
\(H_b(p)\defeq\sum_{i:\,p_i>0}p_i\log_b(1/p_i)\). Base \(2\) gives
\emph{bits}, base \(e\) gives \emph{natural units} (nats), and base
\(10\) gives decimal digits \cite[pp.~1--2]{Shannon1948MathematicalTheory}.
If \(x\) is a chance variable taking value \(i\) with probability
\(p_i\), Shannon writes \(H(x)\) for \(H(p)\); there \(x\) labels the
variable whose law is used and is not an argument of a function
\cite[p.~11]{Shannon1948MathematicalTheory}.
\index{entropy!discrete}\index{entropy!units}\index{bit}
\end{hbdefinition}

Shannon defines \(H=-K\sum_{i=1}^np_i\log p_i\) for finitely many
probabilities and says that the constant \(K\) merely fixes a unit of
measure \cite[p.~11]{Shannon1948MathematicalTheory}. The definition
above keeps his constant and adds two things the formula needs on its
actual domain: the treatment of zero probabilities and the countable
case, in which the series may diverge. The convention \(0\log0=0\) is
also the continuous extension of \(t\log t\) at \(t=0\), proved in
\cref{thm:p15-w02b-p15-040-domain}.

\begin{hbtheorem}{thm:p15-w02b-p15-040-domain}{proved=annex:p15.w02b-p15-040.domain}{Well-definedness, units and extreme values}
Under \cref{def:p15-w02b-p15-040-entropy}:
\begin{enumerate}
\item Every term \(p_i\log(1/p_i)\) lies in \([0,1/e]\). Hence \(H(p)\)
is a well-defined element of \([0,+\infty]\), independent of the order
of summation, and unchanged by adding or deleting indices of
probability zero or by relabelling \(I\) through a bijection.
\item The function \(\varphi(t)\defeq-t\log t\) for \(0<t\leq1\),
\(\varphi(0)\defeq0\), is continuous on \([0,1]\). Consequently
\(H\) is a finite continuous function on \(\Delta_n\) for each
\(n\geq1\).
\item For every base \(b>1\), \(H_b(p)=H_e(p)/\log b\); in particular
\(H_2=H_e/\log2\approx1.4427\,H_e\). The functional with constant
\(K\) equals \(H_b\) for \(b=e^{1/K}\), and \(H_b\) is the functional
with constant \(K=1/\log b\). Choosing \(K\) and choosing a base are the
same choice.
\item \(H(p)=0\) if and only if \(p_i=1\) for one index \(i\). Otherwise
\(H(p)>0\).
\item For \(p\in\Delta_n\), \(H(p)\leq K\log n\), with equality if and
only if \(p_i=1/n\) for every \(i\).
\item On a countably infinite \(I\) the functional is unbounded, and it
takes the value \(+\infty\). For instance, with
\(I=\{2,3,\ldots\}\),
\[
 p_k\defeq\frac{c}{k(\log k)^2},\qquad
 c\defeq\Bigl(\sum_{k\geq2}\frac1{k(\log k)^2}\Bigr)^{-1},
\]
\(c\) is a positive real number, \(p\in\Delta(I)\), and
\(H(p)=+\infty\).
\end{enumerate}
\end{hbtheorem}

Items~4 and~5 are Shannon's first two properties of \(H\): it vanishes
only when the outcome is certain, and for a given number of outcomes it
is largest when they are equally likely
\cite[p.~11]{Shannon1948MathematicalTheory}. Item~5 restates, for an
arbitrary unit, the bound \(0\leq H\leq\log m\) of
\cref{thm:p15-finite-alphabet-entropy} and the uniform maximum of
\cref{thm:p15-rat41-indifference}. Item~6 shows that the
\(K\log n\) bound has no countable analogue. A probability vector with
infinite entropy is a legitimate distribution. Results that subtract
entropies must therefore assume them finite.

\subsection{Shannon's three requirements}
\label{sec:p15-w02b-p15-040-requirements}

Shannon asks what a measure \(H(p_1,\ldots,p_n)\) of the choice
involved in selecting one of \(n\) events of known probabilities should
satisfy \cite[p.~10]{Shannon1948MathematicalTheory}. His third
requirement is given by an example: a choice among probabilities
\(\tfrac12,\tfrac13,\tfrac16\) can be made by first choosing between
two halves and then, half the time, choosing between
\(\tfrac23\) and \(\tfrac13\), so
\[
 H\bigl(\tfrac12,\tfrac13,\tfrac16\bigr)
   =H\bigl(\tfrac12,\tfrac12\bigr)+\tfrac12H\bigl(\tfrac23,\tfrac13\bigr).
\]
The coefficient \(\tfrac12\) is the probability that the second choice
occurs. The definition below states the requirements as properties of a
family of functions, one for each number of outcomes.

\begin{hbdefinition}{def:p15-w02b-p15-040-requirements}{Continuity, monotonicity and grouping}
Let \(H_n:\Delta_n\to\mathbb R\) be given for every \(n\geq1\), and put
\(A(n)\defeq H_n(\tfrac1n,\ldots,\tfrac1n)\). Consider:
\begin{description}
\item[(S1)] each \(H_n\) is continuous on \(\Delta_n\subseteq\mathbb R^n\);
\item[(S2)] \(A(n)<A(n+1)\) for every \(n\geq1\);
\item[(S3)] for every \(n\geq1\), \(p\in\Delta_n\), integers
\(m_1,\ldots,m_n\geq1\) and \(q^{(i)}\in\Delta_{m_i}\),
\[
 \begin{aligned}
 &H_M\bigl(p_1q^{(1)}_1,\ldots,p_1q^{(1)}_{m_1},\;\ldots,\\
 &\qquad p_nq^{(n)}_1,\ldots,p_nq^{(n)}_{m_n}\bigr)\\
 &\qquad=H_n(p)+\sum_{i=1}^np_iH_{m_i}\bigl(q^{(i)}\bigr),
 \end{aligned}
\]
where \(M=m_1+\cdots+m_n\) and the outcomes are listed group by group.
\end{description}
Write (S3u) for the instances of (S3) in which every \(q^{(i)}\) is the
uniform vector on \(m_i\) outcomes.
\index{entropy!grouping property}
\end{hbdefinition}

Shannon's displayed example is the instance of (S3) with \(n=2\),
\(p=(\tfrac12,\tfrac12)\), \(m_1=1\), \(m_2=2\) and
\(q^{(2)}=(\tfrac23,\tfrac13)\). His first two requirements are (S1)
and (S2); he states the second as \(A\) being a monotonic increasing
function of \(n\), meaning that equally likely outcomes give more choice
when there are more of them \cite[p.~10]{Shannon1948MathematicalTheory}.

\begin{hbtheorem}{thm:p15-w02b-p15-040-characterization}{proved=annex:p15.w02b-p15-040.characterization}{Shannon's Theorem~2}
\begin{enumerate}
\item If a family \((H_n)_{n\geq1}\) satisfies (S1), (S2) and (S3u),
then there is a constant \(K>0\) such that
\[
 H_n(p_1,\ldots,p_n)=-K\sum_{i=1}^np_i\log p_i
 \qquad(n\geq1,\ p\in\Delta_n),
\]
with \(0\log0=0\). The constant is \(K=A(2)/\log2\).
\item Conversely, for every \(K>0\) the family
\(H_n(p)=-K\sum_ip_i\log p_i\) satisfies (S1), (S2) and the full
grouping requirement (S3).
\item If (S2) is weakened to \(A(n)\leq A(n+1)\), then (S1) and (S3u)
force the same formula with some \(K\geq0\); the case \(K=0\), in
which every \(H_n\) is identically zero, is then allowed.
\end{enumerate}
\end{hbtheorem}

Shannon states that the only \(H\) satisfying his three assumptions has
the form \(-K\sum_{i=1}^np_i\log p_i\) with \(K\) a positive constant,
and proves it in Appendix~2 \cite[pp.~11, 28--29]{Shannon1948MathematicalTheory}.
The annex writes out his argument in full. It first shows
\(A(s^m)=mA(s)\) by splitting a choice among \(s^m\) equally likely
outcomes into \(m\) successive choices among \(s\). Monotonicity then
pins \(A(t)/A(s)\) between consecutive multiples of \(1/n\), forcing
\(A(t)=K\log t\). Next, a choice among \(\sum n_i\) equally likely
outcomes is split into a choice among \(n\) groups of sizes \(n_i\);
this gives the formula at every vector of rational probabilities.
Continuity extends it to all vectors. Two points that the source leaves
implicit are made explicit. Only the uniform instances (S3u) of the
grouping requirement are used; the general instances (S3) are then a
consequence (item~2). Strict monotonicity is what makes \(K\) positive
(item~3).

\begin{hbwarning}{def:p15-w02b-p15-040-requirements-limits}{What the characterization does and does not justify}
Shannon says that this theorem and the assumptions needed for its proof
are not necessary for the rest of his theory. In his account the theorem
gives the later definitions a certain plausibility; their real
justification lies in their implications, such as the coding theorems
\cite[p.~11]{Shannon1948MathematicalTheory}. The theorem therefore
answers a narrow question: \emph{if} uncertainty about one outcome is
to be measured by a number with properties (S1)--(S3), that number is
\(-K\sum p_i\log p_i\). It does not show that every reasonable
measure of uncertainty must decompose over successive choices.

The grouping requirement does the work. For instance the function
\(G_n(p)\defeq-\log\sum_{i=1}^np_i^2\) is continuous on \(\Delta_n\),
because \(\sum_ip_i^2\geq\bigl(\sum_ip_i\bigr)^2/n=1/n>0\) there by the
Cauchy--Schwarz inequality, and it has
\(G_n(\tfrac1n,\ldots,\tfrac1n)=\log n\). So it satisfies (S1) and
(S2). On Shannon's own example it fails (S3):
\[
 \begin{aligned}
 G\bigl(\tfrac12,\tfrac13,\tfrac16\bigr)&=\log\tfrac{18}7\approx0.9445,\\
 G\bigl(\tfrac12,\tfrac12\bigr)+\tfrac12G\bigl(\tfrac23,\tfrac13\bigr)
   &=\log2+\tfrac12\log\tfrac95\\
   &\approx0.9870 .
 \end{aligned}
\]
Here \(\sum p_i^2\) equals \(\tfrac14+\tfrac19+\tfrac1{36}=\tfrac7{18}\),
\(\tfrac12\), and \(\tfrac49+\tfrac19=\tfrac59\) respectively, in
natural units. By item~2 of
\cref{thm:p15-w02b-p15-040-characterization} every positive multiple of
\(-\sum p_i\log p_i\) satisfies (S3), so \(G\) is not such a multiple; it
is excluded only because it does not decompose over successive choices.
\end{hbwarning}

\subsection{What the number measures}
\label{sec:p15-w02b-p15-040-meaning}

Shannon separates meaning from the engineering problem. Messages often
refer to things, he notes, but for communication the significant fact
is that the actual message is one selected from a set of possible
messages \cite[p.~1]{Shannon1948MathematicalTheory}. Entropy inherits
this standing. It is a functional of the probability vector alone, and
by \cref{thm:p15-w02b-p15-040-domain}(1) it does not change when the
outcomes are relabelled. It is therefore blind to what the outcomes
are, to which of them is true, and to what anyone gains or loses from
them.

\begin{hbwarning}{def:p15-w02b-p15-040-not-truth}{Amount of information is not truth, support or value}
Let \(p\) be a probability vector on the outcomes of a question.
\begin{enumerate}
\item \(H(p)\) does not depend on which outcome is actual. A lower
entropy means a more concentrated distribution, not a more accurate one.
\item A change in \(p\) produced by evidence can raise or lower \(H\)
whether that evidence raises or lowers the probability of a given
hypothesis. Entropy reduction is neither necessary nor sufficient for
evidential support in the sense of an increase in the probability of
the hypothesis.
\item Two distributions with the same entropy can recommend opposite
actions, and the decision value of resolving a question depends on the
available acts and their utilities, which \(H\) does not take as input.
Expected utility is treated in \cref{sec:p17-finite-expected-utility}.
\end{enumerate}
\end{hbwarning}

\begin{hbexample}{def:p15-w02b-p15-040-contrast}{Three computed contrasts}
Use base \(2\); write \(h(t)\defeq-t\log_2t-(1-t)\log_2(1-t)\) for
the entropy of a two-point vector \((t,1-t)\). The values used are
\(h(\tfrac12)=1\), \(h(0.1)=h(0.9)\approx0.4690\) and
\(h(0.01)=h(0.99)\approx0.0808\).

\emph{Truth.} Two forecasts for the outcomes rain and dry are
\(F_1=(0.5,0.5)\) and \(F_2=(0.99,0.01)\). Then \(H(F_1)=1\) bit and
\(H(F_2)\approx0.0808\) bit. Suppose the day is dry. The probability
given to the actual outcome is \(0.5\) under \(F_1\) and \(0.01\) under
\(F_2\); its surprisal \(\log_2(1/0.01)\approx6.644\) bits under \(F_2\)
exceeds the \(1\) bit under \(F_1\). Neither entropy changed when the
outcome was revealed. The forecast \((0.01,0.99)\), which says the
opposite of \(F_2\), has the same entropy \(0.0808\).

\emph{Evidential support.} Let a hypothesis \(h_0\) have prior
probability \(0.1\), and let evidence \(e\) have probability \(0.9\)
under \(h_0\) and \(0.1\) under its negation. Bayes's rule gives
\[
 P(h_0\mid e)=\frac{0.1\cdot0.9}{0.1\cdot0.9+0.9\cdot0.1}=0.5 .
\]
The evidence raises the probability of \(h_0\) fivefold, yet the entropy
of the distribution over \(\{h_0,\neg h_0\}\) rises from \(0.4690\) to
\(1\) bit. Reversing the roles (prior \(0.5\), posterior \(0.1\))
lowers the entropy from \(1\) to \(0.4690\) while the evidence
disconfirms \(h_0\).

\emph{Decision value.} A ticket pays \(1\) if the first of two outcomes
occurs and \(0\) otherwise. The distributions \((0.9,0.1)\) and
\((0.1,0.9)\) have the same entropy \(0.4690\) bits, but the ticket's
expected payoff is \(0.9\) under the first and \(0.1\) under the second.
At a price of \(0.5\), an agent who maximizes expected payoff buys under
the first and declines under the second. Conversely, if every available act has the same
payoff in both outcomes, learning the outcome has no decision value,
whatever the entropy of the distribution over them.
\end{hbexample}

\subsection{Worked dependence and coordinate calculations}
\label{sec:p15-w02b-p15-040-worked}

The joint, marginal and conditional entropies of two finite variables,
their identities and inequalities, and the entropy of a density are
defined and proved in \cref{sec:p15-w02b-p15-041}: see
\cref{def:p15-w02b-p15-041-joint}, \cref{thm:p15-w02b-p15-041-subadditivity},
\cref{thm:p15-w02b-p15-041-chain}, \cref{thm:p15-w02b-p15-041-averaging},
\cref{def:p15-w02b-p15-041-density} and
\cref{thm:p15-w02b-p15-041-coordinates}. The two calculations below
apply those definitions to a dependence comparison and to a change of
coordinates.

\begin{hbexample}{def:p15-w02b-p15-040-dependence}{Same marginals, different dependence}
Use base \(2\). Let \(x\in\{1,2\}\) and \(y\in\{1,2\}\). Compare two
joint tables, rows indexed by \(x\) and columns by \(y\):
\[
 \text{(I)}\quad
 \begin{pmatrix}\tfrac18&\tfrac38\\[2pt]\tfrac18&\tfrac38\end{pmatrix},
 \qquad
 \text{(D)}\quad
 \begin{pmatrix}\tfrac14&\tfrac14\\[2pt]0&\tfrac12\end{pmatrix}.
\]
Both have marginals \((\tfrac12,\tfrac12)\) for \(x\) and
\((\tfrac14,\tfrac34)\) for \(y\). Hence in both
\[
 H(x)=1,\qquad H(y)=h(\tfrac14)=2-\tfrac34\log_23\approx0.8113,
\]
\[
 H(x)+H(y)\approx1.8113 .
\]
In table (I) every cell is the product of its marginals, so
\(H(x,y)=H(x)+H(y)\approx1.8113\). Each row is the law
\((\tfrac14,\tfrac34)\), so \(H_x(y)=H(y)\approx0.8113\).

In table (D), the three positive cells give
\[
 H(x,y)=\tfrac14\log_24+\tfrac14\log_24+\tfrac12\log_22=1.5 .
\]
Given \(x=1\) the law of \(y\) is \((\tfrac12,\tfrac12)\), with
entropy \(1\); given \(x=2\) it is \((0,1)\), with entropy \(0\).
Averaging with weights \(\tfrac12,\tfrac12\) gives \(H_x(y)=0.5\), and
indeed \(H(x)+H_x(y)=1.5=H(x,y)\). Now \(H(x,y)<H(x)+H(y)\), with shortfall
\(H(x)+H(y)-H(x,y)=\tfrac32-\tfrac34\log_23\approx0.3113\) bits, and
\(H_x(y)=0.5<0.8113\approx H(y)\).
In the other direction \(H_y(x)=H(x,y)-H(y)\approx0.6887\).
The marginal entropies alone cannot tell (I) from (D); only the joint
table records the dependence.
\end{hbexample}

\begin{hbexample}{def:p15-w02b-p15-040-coordinates}{A change of coordinates computed both ways}
Use natural units. Let \(x\) be uniform on \((0,1)\), with density
\(1\) there, so \(H(x)=-\int_0^1 1\cdot\log1\,dx=0\). Put
\(y=-\log x\), a strictly decreasing continuously differentiable map of
\((0,1)\) onto \((0,\infty)\).

\emph{Directly.} For \(t>0\),
\(P(y>t)=P(x<e^{-t})=e^{-t}\), so \(y\) has density \(e^{-t}\) on
\((0,\infty)\), the exponential law of rate \(1\). Its entropy is
\[
 H(y)=-\int_0^\infty e^{-t}\log e^{-t}\,dt=\int_0^\infty te^{-t}\,dt=1,
\]
the mean of that law (\cref{thm:p15-gamma-normalization}).

\emph{Through the Jacobian.} Here \(dy/dx=-1/x\), so Shannon's Jacobian
of the old coordinate with respect to the new is
\(J(x/y)=dx/dy=-x\), with \(|J(x/y)|=x\). The expected logarithm is
\[
 \int_0^1\log x\,dx=\bigl[x\log x-x\bigr]_{0}^{1}=-1,
\]
using \(x\log x\to0\) as \(x\to0\). Hence
\(H(x)-\int_0^1 p(x)\log|J(x/y)|\,dx=0-(-1)=1=H(y)\), as
\cref{thm:p15-w02b-p15-041-coordinates} requires. Both integrability
conditions of that result hold: \(\int_0^1|\log x|\,dx=1<\infty\).

The two variables determine each other, yet their entropies are
\(0\) and \(1\). For a finite distribution no such change occurs: a
one-to-one relabelling of outcomes leaves \(H\) unchanged by
\cref{thm:p15-w02b-p15-040-domain}(1). The entropy of a density is
measured against Lebesgue measure in the stated coordinate; it is not a
property of the distribution of outcomes alone.
\end{hbexample}

\subsection{Procedure, limits and relationships}
\label{sec:p15-w02b-p15-040-procedure}

\begin{hbmethod}{def:p15-w02b-p15-040-method}{Computing and reporting a discrete entropy}
\begin{enumerate}
\item Fix the outcome set \(I\) and check that it is finite or countable
and that the probabilities are nonnegative and sum to one.
\item Fix the unit: a base \(b>1\), or the constant \(K=1/\log b\).
Report the unit with the number.
\item Sum \(p_i\log_b(1/p_i)\) over the outcomes with \(p_i>0\). On a
countable set, decide whether the series converges; if it diverges the
entropy is \(+\infty\) and no difference of entropies involving it is
defined.
\item To convert, divide a value in natural units by \(\log b\).
\item For two or more variables, start from the joint table and use the
definitions and identities of \cref{sec:p15-w02b-p15-041}; marginal
entropies do not determine the joint entropy
(\cref{def:p15-w02b-p15-040-dependence}).
\item For a density, state the coordinate and the reference measure, and
check the integrability assumptions of
\cref{def:p15-w02b-p15-041-density}; report how the value changes under
the coordinate changes that are relevant
(\cref{def:p15-w02b-p15-040-coordinates}).
\item Read the result as a property of the stated distribution only: not
as the truth of an outcome, the support for a hypothesis, or the value of
a decision (\cref{def:p15-w02b-p15-040-not-truth}).
\end{enumerate}
\end{hbmethod}

\emph{Limits.} The characterization is a theorem about functions of
probability vectors with properties (S1)--(S3); it says nothing about
how the probabilities were obtained or whether they are correct.
Finite entropy and the bound \(K\log n\) belong to finite outcome sets;
on countable sets the entropy may be infinite. The entropy of a density
is relative to the coordinate system and can be negative.

\emph{Relationships.} The finite-alphabet entropy and its asymptotic
equipartition property in \cref{sec:p15-finite-alphabet-information} use
this functional in natural units, and give it its coding meaning for
long independent sequences. The uniform maximum of
\cref{thm:p15-rat41-indifference} is item~5 of
\cref{thm:p15-w02b-p15-040-domain} in base \(2\). Relative entropy and
its nonnegativity appear in \cref{thm:p15-rmdl-gibbs}. Maximum-entropy
selection of a distribution under constraints, and its dependence on the
reference measure, is treated in
\cref{sec:p15-w02b-p15-044-maxent-existence}.

\emph{Sources.} Shannon's corrected reprint of the 1948 paper:
Introduction, on meaning and logarithmic units
\cite[pp.~1--2]{Shannon1948MathematicalTheory}; Section~6, the three
requirements, Theorem~2 and properties 1 and~2
\cite[pp.~10--11]{Shannon1948MathematicalTheory}; Appendix~2, the
derivation of the logarithmic form
\cite[pp.~28--29]{Shannon1948MathematicalTheory}.
